% !TEX root = ../main.tex

\chapter{关键 Prompt 示例}

\section{项目仓库说明}

为便于复现本文系统的主要实现流程，项目源代码已托管在 GitHub 仓库 \url{https://github.com/phcool/MarketMind.git}。该仓库包含在线展示系统、数据采集与预处理流程、监督微调与强化学习训练实现，以及离线评测与实验分析所需的主要代码。

本附录给出本文实验中使用的主要 Prompt 模板。为保证版面可读性，以下内容对部分重复字段进行了省略，但保留任务说明、输入组织方式和输出约束。

\section{7 日 K 线单周期预测 Prompt}

该 Prompt 用于根据过去 7 个交易日的归一化行情特征预测第 8 个交易日的涨跌方向。输入字段包括开盘价、最高价、最低价、收盘价、成交量、振幅、涨跌幅和换手率。

\begin{quote}\small
过去 7 个交易日的 K 线数据如下（归一化后）：

Day1: open=0.4183, high=0.3760, low=0.3776, close=0.2944, volume=0.6060, amplitude=3.25, pct\_change=-2.2840, turnover=0.3349

Day2: open=0.3028, high=0.2840, low=0.3651, close=0.2742, volume=0.4130, amplitude=1.57, pct\_change=-0.4307, turnover=0.1767

……

Day7: open=0.2589, high=0.2240, low=0.3154, close=0.2419, volume=0.1768, amplitude=1.21, pct\_change=-0.8794, turnover=0.0527

请根据上文 Day1--Day7 的归一化 K 线，只预测下一个交易日（Day8）相对前一日是涨还是跌。

【输出要求（必须严格遵守）】只输出一个字：“涨”或“跌”，不要输出其他任何文字、数字、标点或换行。
\end{quote}

\section{监督微调 CoT 样本 Prompt}

在监督微调数据中，输入仍为 7 日归一化 K 线，但目标输出由单字标签扩展为“分析过程 + 最终标签”。该设计用于让模型学习更稳定的任务解释方式和最终答案格式。

\begin{quote}\small
过去 7 个交易日的 K 线数据如下（归一化后）：

Day1: open=0.4183, high=0.3760, low=0.3776, close=0.2944, volume=0.6060, amplitude=3.25, pct\_change=-2.2840, turnover=0.3349

Day2: open=0.3028, high=0.2840, low=0.3651, close=0.2742, volume=0.4130, amplitude=1.57, pct\_change=-0.4307, turnover=0.1767

……

Day7: open=0.2589, high=0.2240, low=0.3154, close=0.2419, volume=0.1768, amplitude=1.21, pct\_change=-0.8794, turnover=0.0527

请根据上文 Day1--Day7 的归一化 K 线，预测下一个交易日（Day8）相对前一日是涨还是跌。

【输出要求（必须严格遵守）】

请先分析 K 线趋势、量能与波动；将推理过程写在下面一对标记之间（只写推理过程）：

【思维链开始】

（在此撰写推理过程）

【思维链结束】

思维链结束后，请单独一行输出最终答案，仅一个字：“涨”或“跌”，不要引号、标点或其他文字。
\end{quote}

\section{多源信息 1/3/7 日预测 Prompt}

多源信息评测 Prompt 在历史 K 线之外加入新闻摘要和研报摘要，并要求模型同时输出未来 1、3、7 个交易日的方向组合。

\begin{quote}\small
下面是一条股票日度样本，请根据该股票在目标日期当天可见的信息，预测未来短期走势。

你会看到：

1. 该股票在目标日期及之前连续 5 个交易日的归一化 K 线；

2. 截止目标日期最近的 news 摘要；

3. 截止目标日期最近的 report 摘要。

目标日期：2026-01-05

股票代码：000001

过去 5 个交易日的归一化 K 线：

Day1 (2025-12-26): open=0.9123, high=0.7982, low=0.9244, close=0.8934, volume=0.0084, amplitude=0.5200, pct\_change=-0.1141, turnover=0.0000

……

截止该日最近的 news 摘要：

1. date=2026-01-05

title=主力资金连续净流入相关报道

summary=新闻摘要内容……

截止该日最近的 report 摘要：

1. date=2025-12-20

title=机构研究报告标题

summary=研报摘要内容……

请综合 K 线形态、量能变化、news 摘要和 report 摘要，分别预测该股票相对当前交易日收盘价在以下三个阶段的涨跌：1 个交易日后、3 个交易日后、7 个交易日后。

【输出要求（必须严格遵守）】请先分析，再将推理过程写在“【思维链开始】”和“【思维链结束】”之间。思维链结束后，请单独一行按顺序输出三个阶段的答案，格式必须严格为：涨，涨，跌。
\end{quote}

\section{仅行情输入对照 Prompt}

仅行情输入对照实验使用与多源 Prompt 相同的一批样本和同一组标签，但移除 news 与 report 摘要，只保留 5 日归一化 K 线，从而用于观察外部文本信息带来的增益。

\chapter{数据样例}

\section{7 日 K 线训练数据格式}

7 日 K 线训练集字段为 \texttt{prompt} 和 \texttt{label}。其中 \texttt{prompt} 是过去 7 个交易日的归一化 K 线文本，\texttt{label} 是第 8 个交易日相对第 7 个交易日的真实涨跌方向。

\begin{table}[htbp]
  \centering
  \caption{7 日 K 线训练样本字段}
  \begin{tabular}{lll}
    \toprule
    字段 & 含义 & 示例 \\
    \midrule
    \texttt{prompt} & 7 日归一化行情文本 & Day1--Day7 K 线序列 \\
    \texttt{label} & 第 8 日涨跌方向 & 涨 / 跌 \\
    \bottomrule
  \end{tabular}
\end{table}

样例中，单日行情行的组织形式如下：

\begin{quote}\small
Day1: open=0.4183, high=0.3760, low=0.3776, close=0.2944, volume=0.6060, amplitude=3.25, pct\_change=-2.2840, turnover=0.3349
\end{quote}

其中 open、high、low、close 采用个股内部窗口上的 min-max 归一化；volume 先取 \(\log(1+x)\) 后再归一化；pct\_change 采用 z-score 标准化；turnover 采用 min-max 归一化；amplitude 保留原始振幅数值。

\section{CoT 监督微调数据格式}

CoT 风格监督微调数据字段为 \texttt{prompt} 和 \texttt{completion}。其中 \texttt{prompt} 对应模型输入，\texttt{completion} 对应模型需要学习生成的分析过程与最终标签。

该数据通过批处理方式生成。系统将基础 7 日 K 线数据集中的 \texttt{prompt,label} 封装为 DashScope Batch 请求，使用 \texttt{qwen3-max} 生成解释性推理；批处理完成后，再把远端返回结果与原始样本合并为 SFT 所需的 \texttt{prompt,completion} 格式。

\begin{table}[htbp]
  \centering
  \caption{CoT 监督微调样本字段}
  \begin{tabular}{lll}
    \toprule
    字段 & 含义 & 示例 \\
    \midrule
    \texttt{prompt} & 7 日 K 线预测问题 & 归一化 K 线 + 输出要求 \\
    \texttt{completion} & 目标回复 & 分析过程 + 最终“涨/跌” \\
    \bottomrule
  \end{tabular}
\end{table}

训练过程中在 tokenize 时使用 Qwen chat template 将样本转为 user-assistant 对话，并对 prompt 部分 label 置为 \(-100\)，使损失只作用于 assistant 回复部分。

\section{多源信息评测数据格式}

多源信息数据字段包括目标日期、股票代码、5 日 K 线、新闻摘要、研报摘要和未来多周期标签。

\begin{table}[htbp]
  \centering
  \caption{多源信息评测样本字段}
  \begin{tabular}{@{}p{0.34\textwidth}p{0.54\textwidth}@{}}
    \toprule
    字段 & 含义 \\
    \midrule
    \texttt{date} & 目标交易日期 \\
    \texttt{stock} & 股票代码 \\
    \texttt{kline\_5d} & 目标日及之前连续 5 个交易日的归一化 K 线 JSON \\
    \texttt{news} & 截至目标日最近 30 日内最多 30 条新闻摘要 JSON \\
    \texttt{reports} & 截至目标日最近 30 日内最多 3 条研报摘要 JSON \\
    \texttt{future labels} & 未来 1、3、7 个交易日方向标签 \\
    \bottomrule
  \end{tabular}
\end{table}

其中，未来标签字段为 \texttt{future\_1\_3\_7\_trade\_day\_labels}。未来标签通过比较目标日期收盘价与未来第 1、第 3、第 7 个交易日收盘价得到。若未来收盘价大于或等于目标日期收盘价，则标记为“涨”，否则标记为“跌”。例如“涨，涨，跌”表示未来 1 日和 3 日方向为上涨，未来 7 日方向为下跌。

\section{多周期 GRPO 训练数据格式}

多周期 GRPO 训练数据同样基于 7 日归一化 K 线构造，但标签不再是单个下一交易日方向，而是未来 1、3、7 个交易日方向组合。该数据由 \texttt{build\_quotes\_7d\_multi\_dataset.py} 生成，训练文件为 \texttt{quotes\_7d\_multi\_pre2026\_dataset.csv}，验证文件为 \texttt{quotes\_7d\_multi\_eval\_20260101\_20260228.csv}。

\begin{table}[htbp]
  \centering
  \caption{多周期 GRPO 训练样本字段}
  \begin{tabular}{@{}p{0.20\textwidth}p{0.68\textwidth}@{}}
    \toprule
    字段 & 含义 \\
    \midrule
    \texttt{prompt} & 过去 7 个交易日的归一化 K 线文本，以及要求模型先分析、后输出三项最终方向的任务说明 \\
    \texttt{labels} & 未来第 1、第 3、第 7 个交易日相对 Day7 收盘价的方向组合，例如“涨，跌，涨” \\
    \bottomrule
  \end{tabular}
\end{table}

其中，表中的 \texttt{labels} 对应 CSV 中的完整字段 \texttt{future\_1\_3\_7\_trade\_day\_labels}。GRPO 训练脚本也兼容 \texttt{prompt,completion} 两列格式，但要求 \texttt{completion} 中能够解析出未来 1、3、7 个交易日的最终“涨/跌”结果。奖励函数会逐项比较三项方向标签，并取平均正确率作为该 completion 的 reward。

\chapter{实验参数配置}

\section{主要数据集概览}

\begin{table}[htbp]
  \centering
  \caption{本文使用的主要数据集概览}
  \begin{tabular}{p{0.23\textwidth}p{0.19\textwidth}p{0.22\textwidth}p{0.16\textwidth}}
    \toprule
    数据集 & 主要输入 & 标签形式 & 样本数 \\
    \midrule
    7 日 K 线监督数据 & 连续 7 个交易日归一化 K 线 & 下一交易日涨/跌 & 339287 \\
    CoT 监督微调数据 & 7 日 K 线 Prompt & 分析过程 + 最终涨/跌 & 599871 \\
    多源信息评测数据 & 5 日 K 线 + 新闻摘要 + 研报摘要 & 未来 1、3、7 日涨跌组合 & 1691 \\
    GRPO 多周期方向数据 & 7 日 K 线 Prompt & 未来 1、3、7 日涨跌组合 & 由行情窗口生成 \\
    \bottomrule
  \end{tabular}
\end{table}

\section{数据集构建参数}

\begin{table}[htbp]
  \centering
  \caption{数据集构建主要参数}
  \begin{tabular}{@{}p{0.23\textwidth}p{0.23\textwidth}p{0.34\textwidth}@{}}
    \toprule
    参数 & 取值 & 说明 \\
    \midrule
    7 日 K 线训练起始日期 & 2021-01-01 & 标签日不早于该日期 \\
    7 日 K 线训练截止日期 & 2026-01-01 & 标签日早于该日期 \\
    K 线输入长度 & 7 个交易日 & 用于单周期涨跌预测 \\
    多源 K 线输入长度 & 5 个交易日 & 用于 1/3/7 日多周期预测 \\
    多源评测日期范围 & 2026-01-01 至 2026-04-01 & 目标日期范围 \\
    新闻最大回看窗口 & 30 日 & 不使用目标日期之后信息 \\
    新闻数量上限 & 30 条 & 每个样本最多保留 30 条摘要 \\
    研报数量上限 & 3 条 & 每个样本最多保留 3 条摘要 \\
    \bottomrule
  \end{tabular}
\end{table}

\section{LoRA 监督微调参数}

本文使用 LoRA 对 Qwen2.5-7B-Instruct 进行监督微调。主要参数如下。

\begin{table}[htbp]
  \centering
  \caption{LoRA 监督微调的主要训练参数}
  \begin{tabular}{p{0.32\textwidth}p{0.42\textwidth}}
    \toprule
    参数项 & 配置值 \\
    \midrule
    基础模型 & Qwen2.5-7B-Instruct \\
    微调方式 & LoRA \\
    LoRA 秩 \(r\) & 64 \\
    LoRA \(\alpha\) & 128 \\
    LoRA dropout & 0.05 \\
    目标模块 & 注意力投影层与前馈网络关键线性层 \\
    训练轮数 & 2 \\
    学习率 & \(2\times10^{-4}\) \\
    学习率调度 & cosine \\
    warmup ratio & 0.03 \\
    单卡 batch size & 1 \\
    梯度累积步数 & 16 \\
    有效 batch size & 128（8 卡） \\
    最大序列长度 & 2048 tokens \\
    精度设置 & bfloat16 \\
    稳定性设置 & gradient checkpointing, safetensors \\
    \bottomrule
  \end{tabular}
\end{table}

\begin{table}[htbp]
  \centering
  \caption{LoRA SFT 主要训练参数}
  \begin{tabular}{@{}p{0.23\textwidth}p{0.28\textwidth}p{0.35\textwidth}@{}}
    \toprule
    参数 & 取值 & 说明 \\
    \midrule
    基础模型 & Qwen2.5-7B-Instruct & 指令模型基座 \\
    训练数据 & CoT 训练集 & 由 7 日 K 线样本扩展得到 \\
    最大序列长度 & 2048 & token 截断长度 \\
    训练轮数 & 2 & \texttt{num\_train\_epochs} \\
    学习率 & \(2\times10^{-4}\) & LoRA adapter 学习率 \\
    batch size & 1 & 单卡 batch size \\
    梯度累积步数 & 16 & 每累积 16 步更新一次参数 \\
    warmup ratio & 0.03 & 预热比例 \\
    学习率调度 & cosine & 余弦退火 \\
    随机种子 & 42 & 实验复现 \\
    精度 & bfloat16 & 训练数值精度 \\
    \bottomrule
  \end{tabular}
\end{table}

\begin{table}[htbp]
  \centering
  \caption{LoRA adapter 参数}
  \begin{tabular}{lll}
    \toprule
    参数 & 取值 & 说明 \\
    \midrule
    LoRA rank & 64 & 低秩矩阵秩 \\
    LoRA alpha & 128 & 缩放系数 \\
    LoRA dropout & 0.05 & dropout 概率 \\
    target modules & q/k/v/o/gate/up/down projection & 注意力与 MLP 投影层 \\
    bias & none & 不训练 bias \\
    task type & CAUSAL\_LM & 自回归语言模型 \\
    \bottomrule
  \end{tabular}
\end{table}

多卡训练默认使用 8 张 GPU，环境变量中设置 \texttt{CUDA\_VISIBLE\_DEVICES=0,1,2,3,4,5,6,7}，并通过 \texttt{torchrun} 启动训练进程。

\section{GRPO 强化学习参数}

多周期方向 GRPO 强化学习阶段使用纯 7 日 K 线构造的未来 1、3、7 个交易日方向数据，其主要训练配置如下，其中包括训练/验证数据来源、组内采样规模、上下文长度限制以及平均方向奖励设置。

\begin{table}[htbp]
  \centering
  \caption{多周期方向 GRPO 训练的主要参数设置}
  \begin{tabular}{p{0.32\textwidth}p{0.42\textwidth}}
    \toprule
    参数项 & 配置值 \\
    \midrule
    基础模型 & Qwen2.5-7B-Instruct \\
    训练数据 & 7 日 K 线多周期方向数据 \\
    验证数据 & 2026 年多周期方向评测集 \\
    学习率 & \(1\times10^{-6}\) \\
    单卡 batch size & 1 \\
    梯度累积步数 & 4 \\
    训练轮数 & 1 \\
    组内采样数 \texttt{num\_generations} & 8 \\
    最大 Prompt 长度 & 4096 tokens \\
    最大 completion 长度 & 2048 tokens \\
    主奖励形式 & 未来 1/3/7 日方向逐项匹配，奖励取三项平均值 \\
    Rollout 方式 & vLLM colocate \\
    启动方式 & 8 卡 \texttt{accelerate launch} \\
    恢复机制 & 自动恢复最近 checkpoint \\
    \bottomrule
  \end{tabular}
\end{table}

\section{vLLM 评测参数}

多源信息评测使用 vLLM 加载基础模型和 LoRA adapter。主要参数如下。

\begin{table}[htbp]
  \centering
  \caption{vLLM 评测主要参数}
  \begin{tabular}{lll}
    \toprule
    参数 & 取值 & 说明 \\
    \midrule
    基础模型 & Qwen/Qwen2.5-7B-Instruct & 推理基座 \\
    LoRA adapter & LoRA 微调权重 & 监督微调 adapter \\
    temperature & 0.0 & 确定性解码 \\
    top\_p & 1.0 & 不额外截断概率质量 \\
    max new tokens & 4096 & 多源输入组最大生成长度 \\
    tensor parallel size & 4 & 张量并行度 \\
    pipeline parallel size & 2 & 流水并行度 \\
    max model len & 32768 & 最大上下文长度 \\
    gpu memory utilization & 0.9 & 显存利用上限 \\
    \bottomrule
  \end{tabular}
\end{table}

评测过程使用正则表达式从模型输出中抽取最终一行“涨，跌，涨”形式的三元标签。如果无法抽取合法标签，则记为“提取失败”。平均准确率按三个预测位置逐项计分。


\section{提取失败样本说明}

在多周期任务中，模型必须按固定格式输出三个方向标签，例如“涨，跌，涨”。如果模型输出包含多余解释、缺少某个预测周期，或没有给出可解析的三元标签，评测过程会将其标记为“提取失败”。这类样本在包含失败统计中按 0 分计入，在排除失败统计中从分母中剔除。

提取失败并不一定表示模型完全没有进行判断，也可能只是输出格式不符合自动解析要求。因此，后续系统可以通过更严格的结构化输出约束降低该问题，例如使用 JSON schema、函数调用式输出或单独的分类头。
